\begin{table}[htb]
  \caption[Moderationsanalyse Religiosität und biologisches Geschlecht auf das psychische Wohlbefinden]{\textit {Moderationsanalyse Religiosität und biologisches Geschlecht auf das psychische Wohlbefinden}} 
  \label{H3_Moderation}
  \centering
  \begin{adjustbox}{width=11cm} %{width=\textwidth}
  \small
  \begin{tabular}{lrrrrrr}
    \hline
               & $B$   & $SE(B)$&  $t$    & $p$   & $\Delta R^{2}$\\
    \hline
  Konstante    & 3.39    & 0.17 & 19.54   & $<$~.001 & \\
  Religiosität & $-$0.15 & 0.18 & $-$0.83 & .406     & \\
  Geschelcht   & $-$0.11 & 0.10 & $-$1.15 & .251     & \\
  Interaktionsterm & 0.22 & 0.10 & 2.15   & .032     & .010* \\
  $R^{2}$          &      &      &        &          & .069** \\
     \hline
  \end{tabular}
  \end{adjustbox}
  
  \begin{tablenotes}
      \item \textit{Anmerkungen.} \(N\)~=~427; Wertebereich der Variable Religiosität -1.69 (\textit{gar nicht religiös}) bis 1.16 (\textit{sehr religiös}); biologisches Geschlecht 1 (\textit{männlich}) bis 2 (\textit{weiblich}).\\ *$p<$~.05, **$p<$~.001
    \end{tablenotes}
  \end{table}


